% Extensión aditiva. Canon leído: CURRENT.json rev. 2026-07-22.2.
% Autores del proyecto: Oumar Haidara Fall y Rubén Ramos Balsa.

\chapter{Droites de grandeur et réalisation dimensionnelle du noyau HMT}
\label{chap:funtor-dimensional-hmt}

La transition de HMT vers la mécanique dimensionnelle ne consiste pas à adjoindre
un symbole d'unité à un nombre réel. Le noyau arithmétique produit des
caractères, des cocycles, des rapports et des ordres d'une carte décimale. Une grandeur
physique, en revanche, est un élément d'une droite pourvue d'un type dimensionnel.
La distinction est structurelle : un changement d'unités modifie les coordonnées
d'une grandeur, mais non la grandeur elle-même. Cette section construit, d'une
part, le réseau ambiant requis pour réaliser indépendamment
masse, longueur, durée et charge, et détermine que son rang minimal est quatre.
D'autre part, elle construit un foncteur monoïdal de trajectoires avec signature qui
prend ses valeurs dans une portion de ce réseau. La minimalité du réseau n'est pas
attribuée au foncteur.

La construction a deux conséquences. En premier lieu, elle permet de réaliser
dans un réseau commun l'action, la durée, la longueur, la charge et la masse
sans introduire le Système international dans les formules génératrices ;
l'assignation de charge aux états requiert, en outre, un caractère entier
spécifique. En second lieu, elle démontre exactement quelle partie d'une échelle
absolue ne peut provenir de scalaires adimensionnels. Cette impossibilité
n'affecte pas les rapports, les caractères de masse ni les identités
constitutives ; elle fixe le type mathématique correct de leurs sorties.
La distinction entre dimension, coordonnée numérique et choix d'unité est celle
de la théorie classique de l'analyse dimensionnelle \cite{Buckingham1914,Bridgman1931} ;
la structure de droites qui suit explicite cette distinction dans le domaine
HMT\@.

\section{Réseau de dimensions et droites orientées}

Soit
\[
 \Lambda_{\rm D}
 =\mathbb Z\langle \mathbf M,\mathbf L,\mathbf T,\mathbf Q\rangle
 \cong\mathbb Z^4
\]
le réseau de dimensions, écrit en termes de masse, de longueur, de durée
et de charge. Pour
\(d=(d_M,d_L,d_T,d_Q)\in\Lambda_{\rm D}\), nous notons
\(\mathcal L_d\) une droite réelle orientée et
\(\mathcal L_d^+\) sa demi-droite positive. Le produit tensoriel réalise la
somme des dimensions,
\[
 \mathcal L_d\otimes\mathcal L_e\simeq\mathcal L_{d+e},
 \qquad
 \mathcal L_d^\vee\simeq\mathcal L_{-d}.
\]
La collection de ces droites, munie des isomorphismes positifs et du produit
tensoriel, forme un groupoïde de Picard symétrique. Un scalaire adimensionnel est
un élément de \(\mathcal L_0\simeq\mathbb R\) ; une grandeur de dimension
\(d\) est un élément de \(\mathcal L_d\).

Nous introduisons quatre dimensions adaptées à HMT :
\[
 \mathbf S=\mathbf M+2\mathbf L-\mathbf T,
 \qquad
 \mathbf C=\mathbf L-\mathbf T,
 \qquad
 \mathbf T=\mathbf T,
 \qquad
 \mathbf Q=\mathbf Q.
\]
Elles représentent, respectivement, l'action, la vitesse, la durée et la charge. Dans la
base \((\mathbf M,\mathbf L,\mathbf T,\mathbf Q)\), leurs colonnes forment la
matrice
\[
 B_{\rm D}=
 \begin{pmatrix}
 1&0&0&0\\
 2&1&0&0\\
 -1&-1&1&0\\
 0&0&0&1
 \end{pmatrix}.
\]

\begin{theorem}[Base dimensionnelle adaptée à HMT]
\label{thm:base-dimensional-hmt}
Les dimensions
\[
 \mathbf S,\qquad\mathbf C,\qquad\mathbf T,\qquad\mathbf Q
\]
constituent une base unimodulaire de
\(\Lambda_{\rm D}\). En particulier, toute dimension mécanique,
électromagnétique ou gravitationnelle engendrée par
\((\mathbf M,\mathbf L,\mathbf T,\mathbf Q)\) possède une expression unique avec des
exposants entiers dans les quatre droites
\(\mathcal L_S,\mathcal L_C,\mathcal L_T,\mathcal L_Q\).
\end{theorem}

\begin{proof}
La matrice \(B_{\rm D}\) est triangulaire par blocs et
\(\det B_{\rm D}=1\). Elle appartient donc à
\(\operatorname{GL}_4(\mathbb Z)\). Son inverse entier fournit
\[
 \mathbf L=\mathbf C+\mathbf T,
 \qquad
 \mathbf M=\mathbf S-2\mathbf C-\mathbf T.
\]
Les quatre dimensions originales sont donc retrouvées de manière unique.
\end{proof}

Parmi les identités que nous utiliserons figurent
\begin{align}
 \mathbf E&=\mathbf S-\mathbf T,
 &\mathbf Z&=\mathbf S-2\mathbf Q,\label{eq:dim-EZ}\\
 \boldsymbol\mu&=\mathbf S-\mathbf C-2\mathbf Q,
 &\boldsymbol\varepsilon&=-\mathbf S-\mathbf C+2\mathbf Q,
 \label{eq:dim-mueps}\\
 \mathbf G&=-\mathbf S+5\mathbf C+2\mathbf T,
 &\boldsymbol\kappa&=-\mathbf S+\mathbf C+2\mathbf T,
 \label{eq:dim-Gkappa}\\
 \boldsymbol\Lambda&=-2\mathbf C-2\mathbf T.
 \label{eq:dim-Lambda}
\end{align}
Ici \(\mathbf E\) est l'énergie, \(\mathbf Z\) l'impédance,
\(\boldsymbol\mu\) la perméabilité, \(\boldsymbol\varepsilon\)
la permittivité, \(\mathbf G\) la dimension de la constante gravitationnelle,
\(\boldsymbol\kappa=\mathbf G-4\mathbf C\), et
\(\boldsymbol\Lambda=-2\mathbf L\).

\section{Groupe de changement d'échelle et théorème d'impossibilité}

Le groupe des changements de base d'unités est
\[
 \mathcal G_{\rm D}
 :=\operatorname{Hom}(\Lambda_{\rm D},\mathbb R_{>0})
 \cong(\mathbb R_{>0})^4.
\]
Pour \(g\in\mathcal G_{\rm D}\) et \(d\in\Lambda_{\rm D}\), nous écrivons
\(g^d\) pour le caractère correspondant. Si \(u_d\) est une base positive
de \(\mathcal L_d\), le changement de carte
\(u_d\mapsto g^d u_d\) transforme la coordonnée d'une grandeur fixe selon
\[
 x\,u_d=(g^{-d}x)(g^du_d).
\]
La grandeur ne change pas ; sa représentation numérique change.

\begin{theorem}[Impossibilité d'une échelle absolue à partir de données
adimensionnelles]
\label{thm:no-go-dimensional-rango-cuatro}
Soit \(X\) un espace sur lequel \(\mathcal G_{\rm D}\) agit
trivialement et soit \(d\ne0\). Toute application
\(F:X\to\mathcal L_d\) naturelle relativement à tous les changements
d'unités et ne recevant pas de base dimensionnelle supplémentaire est identiquement
nulle. En revanche, tout quotient de deux sections non nulles d'une même droite,
et tout produit dont la dimension totale est nulle, est invariant par
\(\mathcal G_{\rm D}\).
\end{theorem}

\begin{proof}
La naturalité et l'action triviale sur \(X\) impliquent
\[
 F(x)=F(gx)=g^dF(x)
 \qquad(g\in\mathcal G_{\rm D}).
\]
Comme \(d\ne0\), il existe \(g\) avec \(g^d\ne1\) ; d'où
\((g^d-1)F(x)=0\), donc \(F(x)=0\). Pour deux sections de dimension
\(d\), les facteurs \(g^d\) s'annulent dans le quotient. La même annulation
se produit dans tout monôme de dimension totale nulle.
\end{proof}

Le théorème n'établit pas qu'une théorie discrète soit incapable de construire une
échelle interne. Il établit que cette échelle doit apparaître comme une section d'une
droite ou, de manière équivalente, que la théorie doit étendre son domaine par le
torseur des repères dimensionnels. C'est précisément cette extension qui définit
la mécanique dimensionnelle.

\begin{corollary}[Rang minimal]
Tout système de droites permettant de réaliser indépendamment la masse,
la longueur, la durée et la charge possède un rang d'échelle au moins égal à quatre. La base
\((\mathbf S,\mathbf C,\mathbf T,\mathbf Q)\) atteint la borne ; elle est donc
minimale.
\end{corollary}

\begin{proof}
Les quatre dimensions originales sont linéairement indépendantes dans
\(\Lambda_{\rm D}\). Trois générateurs ne peuvent engendrer un réseau de
rang quatre. Le \cref{thm:base-dimensional-hmt} donne quatre générateurs qui,
eux, l'engendrent.
\end{proof}

\section{Espace dimensionnel interne : degrés, cartes et opérateurs de transduction HMT–MD}

Un repère interne HMT--MD est un quadruplet de bases positives
\[
 \mathfrak u=(u_S,u_C,u_T,u_Q)
 \in\operatorname{Fr}_{\rm D},
\]
où \(\operatorname{Fr}_{\rm D}\) est un torseur principal de
\(\mathcal G_{\rm D}\). Les bases dérivées sont définies tensoriellement :
\begin{align*}
 u_L&=u_Cu_T,&
 u_M&=u_Su_C^{-2}u_T^{-1},&
 u_E&=u_Su_T^{-1},\\
 u_Z&=u_Su_Q^{-2},&
 u_\mu&=u_Su_C^{-1}u_Q^{-2},&
 u_\varepsilon&=u_S^{-1}u_C^{-1}u_Q^2,\\
 u_G&=u_S^{-1}u_C^5u_T^2,&
 u_\kappa&=u_S^{-1}u_Cu_T^2,&
 u_\Lambda&=u_C^{-2}u_T^{-2}.
\end{align*}
La juxtaposition désigne ici le produit tensoriel de bases. Une carte externe
est un autre repère \(\mathfrak u_{\rm ext}\) muni d'un isomorphisme positif
entre les deux systèmes de droites. Rien n'impose que cette carte soit SI ; le
choix de SI relève de la publication métrologique d'une grandeur.

\subsection{Droite d'action}

Le lecteur déterminantal local part de
\[
 G_H=\operatorname{coker}(H_4),
 \qquad |G_H|=16,
 \qquad \Sigma_H=\varphi\alpha_{\rm HMT}^{16},
\]
et démontre
\[
 10^{-34}<\Sigma_H<10^{-33},
 \qquad \operatorname{ord}_{10}(\Sigma_H)=-34.
\]
La réalisation régionale postérieure au retour fournit la coordonnée
angulaire adimensionnelle
\[
 \etaRet
 =\frac{\Cdeg}{2}-\alpha_{\angle,{\rm deg}}
 =1.0545718169150390611336905847242997\ldots.
\]
Nous définissons la section interne d'action réduite par
\begin{equation}
 \hbar_{\rm int}
 :=\etaRet\,10^{-34}u_S,
 \qquad
 h_{\rm int}:=2\pi\hbar_{\rm int}.
 \label{eq:accion-interna}
\end{equation}
L'exposant est une sortie du lecteur ; \(u_S\) est la base d'une droite.
La composition est covariante sous \(\mathcal G_{\rm D}\) et ne contient pas
d'identification à \(\mathrm{J\,s}\).

\subsection{Périodicité des arêtes de voie et vitesse constitutive résultante}

Soit \(\mathsf P_{\rm TPK}\) la catégorie libre des trajectoires du graphe des
états enrichis. La graduation par longueur
\[
 \tau(\gamma)=|\gamma|,
 \qquad
 \tau(\gamma\eta)=\tau(\gamma)+\tau(\eta)
\]
définit l'horloge combinatoire du protocole. Sa réalisation est
\[
 T_{\rm int}(\gamma)=\tau(\gamma)u_T.
\]
Cette durée compte les transitions du graphe fixé. Une subdivision du graphe
définit une autre horloge et doit être accompagnée de son morphisme de raffinement.

Les canaux \(q_\pm\) engendrent les fonctions symétrique et antisymétrique
\[
 \mathcal E=
 \log\!\bigl[(1-q_+^{90})(1-q_-^{90})\bigr]
 -\log\!\bigl[(1-q_+^{120})(1-q_-^{120})\bigr],
\]
\[
 \mathcal B=
 \log\!\frac{1-q_-^{90}}{1-q_+^{90}}
 -\log\!\frac{1-q_-^{120}}{1-q_+^{120}}.
\]
Celles-ci déterminent
les coefficients positifs
\[
 \widehat c=e^{-\mathcal E},
 \qquad \widehat Z=e^{\mathcal B}.
\]
Leurs réalisations internes sont
\begin{equation}
 c_{\rm int}=\widehat c\,u_C,
 \qquad Z_{\rm int}=\widehat Z\,u_Z.
 \label{eq:cZ-internos}
\end{equation}
Pour un canal constant le long d'une trajectoire, on obtient la longueur
\begin{equation}
 L_{\rm int}(\gamma)
 =\widehat c\,|\gamma|\,u_L.
 \label{eq:longitud-ruta}
\end{equation}
Si le canal dépend de l'arête, la formule covariante est la somme
\(\sum_{e\in\gamma}\widehat c(e)u_L\) ; son additivité par concaténation est
conservée.

\subsection{Échelle élémentaire de charge et caractère conditionnel}

Puisque \(u_Z=u_Su_Q^{-2}\), le quotient
\(h_{\rm int}/Z_{\rm int}\) appartient à \(\mathcal L_Q^{\otimes2}\).
L'orientation positive des droites sélectionne une unique racine positive.
Nous définissons
\begin{equation}
 e_{\rm int}
 :=\left(\frac{2\alpha_{\rm HMT}h_{\rm int}}
 {Z_{\rm int}}\right)^{1/2}\in\mathcal L_Q^+.
 \label{eq:carga-interna}
\end{equation}

\begin{proposition}[Clôture interne de la relation de structure fine]
La section \eqref{eq:carga-interna} est l'unique section positive qui
satisfait
\[
 \alpha_{\rm HMT}
 =\frac{Z_{\rm int}e_{\rm int}^{\,2}}{2h_{\rm int}}.
\]
L'égalité est invariante sous le groupe complet de changement d'échelle.
\end{proposition}

\begin{proof}
L'existence résulte de la positivité de
\(\alpha_{\rm HMT},h_{\rm int},Z_{\rm int}\). Une droite réelle orientée
possède une seule racine positive d'une section positive de son carré. La
dimension du membre de droite est
\(\mathbf Z+2\mathbf Q-\mathbf S=0\), de sorte que tous les facteurs de
changement d'échelle s'annulent.
\end{proof}

Cette proposition ne compte pas comme une seconde détermination indépendante de
\(\alpha\) : elle réalise, au sein du système de droites, la compatibilité
constitutive entre action, impédance et une échelle élémentaire de charge. Elle ne
construit pas de charge sur les états. Pour cela, il faut en outre un
caractère entier
\[
 n_Q:\operatorname{Ob}(\mathsf P_{\rm TPK}^{\rm sig})\longrightarrow
 \mathbb Z;
 \qquad
 Q(x)=n_Q(x)e_{\rm int}.
\]
L'existence et la sélection de \(n_Q\) sont des hypothèses supplémentaires. En
particulier, \(n_Q\) n'est pas une composante du foncteur de trajectoires
construit ci-après. L'échelle élémentaire et la classification des états
sont des applications distinctes.

\section{Foncteur monoïdal de trajectoires avec signature}

Soit \(\mathsf P_{\rm TPK}^{\rm sig}\) la catégorie des trajectoires
enrichies décorées par une coordonnée de signature
\(v_x\in\mathbb Z^5\) en chaque objet. Pour
\(\gamma:x\to y\), nous posons
\[
 \Delta v(\gamma)=v_y-v_x.
\]
La composition satisfait
\(\Delta v(\gamma\eta)=\Delta v(\gamma)+\Delta v(\eta)\).
Soit
\[
 \chi_{\rm mass}(v)=\exp\langle\lambda_{\rm mass},v\rangle
\]
le caractère positif de masse. Sur le quotient des voies, nous fixons la loi
d'action déterminée par cinq données : la trace normalisée des émissions
distinctes, un cocycle positif et additif, le défaut de changement
\(\Delta_4P_3\), la capacité résiduelle \(\lambda=1-\delta\) et la somme locale
des coûts d'arête. Dans cette loi,
\[
 \mathcal A_*(\gamma)
 =N_{\rm giro}(\gamma)+\lambda_*N_{\rm hoja}(\gamma),
 \qquad
 \lambda_*=1-\frac5{468}-\frac3{11}\Delta_4.
\]
L'argument suivant requiert seulement l'additivité ; il s'applique également à la
fonctionnelle \(I_\lambda\) avec tout \(\lambda>0\).

Considérons le monoïde
\[
 \mathfrak M_{\rm D}
 =\mathbb N\times\mathbb R_{\ge0}\times
  \mathbb R_{\ge0}\times\mathbb R_{>0}
\]
avec produit
\[
 (t,\ell,s,r)\odot(t',\ell',s',r')
 =(t+t',\ell+\ell',s+s',rr').
\]

\begin{theorem}[Foncteur monoïdal de trajectoires avec signature]
\label{thm:funtor-dimensional-minimo}
L'application
\[
 \mathfrak D_{\rm HMT}(\gamma)
 =\bigl(
 |\gamma|,
 \widehat c|\gamma|,
 \mathcal A_*(\gamma),
 \chi_{\rm mass}(\Delta v(\gamma))
 \bigr)
\]
définit un foncteur monoïdal
\[
 \mathfrak D_{\rm HMT}:
 \mathsf P_{\rm TPK}^{\rm sig}\longrightarrow\mathfrak M_{\rm D}.
\]
Après choix d'un repère interne, ses trois composantes additives sont réalisées
comme
\[
 T_{\rm int}(\gamma)=|\gamma|u_T,
 \quad
 L_{\rm int}(\gamma)=\widehat c|\gamma|u_L,
 \quad
 S_{\rm int}(\gamma)=\mathcal A_*(\gamma)\hbar_{\rm int},
\]
et la quatrième transporte des rapports adimensionnels de masse. Le foncteur ne contient pas
de coordonnée de charge et n'engendre pas le réseau dimensionnel ambiant.
\end{theorem}

\begin{proof}
La longueur d'une trajectoire, le nombre de rotations et le nombre de changements de
feuillet sont additifs sous concaténation compatible. Par hypothèse,
\(\mathcal A_*\) est additive. La différence des signatures est additive et le caractère
de masse transforme les sommes en produits. Chaque composante respecte donc
la loi \(\odot\), et la voie vide est envoyée sur
\((0,0,0,1)\). La réalisation linéaire conserve les sommes parce qu'elle multiplie
par des sections fixes des droites correspondantes.
\end{proof}

\begin{remark}[Réseau ambiant et portée du foncteur]
Le corollaire de rang minimal se rapporte à un système capable de faire varier
indépendamment masse, longueur, durée et charge. Le foncteur précédent
réalise durée, longueur et action, et transporte un rapport de masse ; il n'utilise pas
\(u_Q\) et ne démontre pas de borne de rang. L'assignation ultérieure de charges
demeure conditionnée à l'existence de \(n_Q\). Dans le réseau dimensionnel,
les types de ses composantes sont
\[
 \mathbf T,\qquad \mathbf L=\mathbf C+\mathbf T,\qquad
 \mathbf S,\qquad \mathbf0,
\]
de sorte que son image dimensionnelle est contenue dans
\(\langle\mathbf S,\mathbf C,\mathbf T\rangle_{\mathbb Z}\), de rang trois.
Le changement d'échelle \(u_Q\mapsto\lambda u_Q\), les autres bases étant fixes, laisse
invariantes toutes les composantes du foncteur ; celui-ci ne réalise pas fidèlement la
quatrième direction ambiante.
\end{remark}

Si l'on assigne à chaque objet la section électronique bidirectionnelle comme
origine de la carte de masse,
\[
 m(x)=m_e^{\beta}\chi_{\rm mass}(v_x-v_e),
\]
alors tout morphisme \(\gamma:x\to y\) satisfait exactement
\[
 \frac{m(y)}{m(x)}
 =\chi_{\rm mass}(\Delta v(\gamma)).
\]
Le théorème ferme la propagation dimensionnelle une fois la signature donnée. Il ne change pas
le statut de l'application précédente qui associe une espèce à une trajectoire et à
sa signature.

\section{Étape basale et clôture bidirectionnelle de l'échelle électronique}

L'invariant électronique en vigueur est le scalaire positif
\[
 \mathfrak e_0=
 \frac{\sqrt3}{4}
 \left(e^{\varphi/\pi^2}-22\alpha_{\rm HMT}^3\right)
 (1+15\Delta_4).
\]
Sa valeur numérique est
\(0.510998922488066\ldots\). C'est l'étape basale de la généalogie, non sa
sortie directrice. Dans le repère dimensionnel interne, il définit une section d'énergie
et une section de masse reliées par la vitesse :
\begin{equation}
 E_e^{\rm int}=\mathfrak e_0u_E,
 \qquad
 m_e^{\rm int}=\frac{E_e^{\rm int}}{c_{\rm int}^{2}}
 =\mathfrak e_0\widehat c^{-2}u_M.
 \label{eq:electron-line-valued}
\end{equation}
La voie centrale et ses deux lecteurs bidirectionnels déterminent, sans consulter
la masse externe,
\[
 \kappa_e=80-54\alpha_{\rm HMT}+6\Delta_4,
 \qquad R_{\rm act}=\frac{H_5}{\etaRet}.
\]
La section électronique en vigueur est donc
\begin{equation}
 \begin{aligned}
 E_e^{\beta}&=E_e^{(0)}R_{\rm act}^{\kappa_e},
 &m_e^{\beta}&=\frac{E_e^{\beta}}{c_{\rm int}^{2}},\\
 E_e^\beta
 &=0.5109989506900008086082571648\ldots\ {\rm MeV},
 &u_E&=1\,{\rm MeV}.
 \end{aligned}
 \label{eq:electron-line-valued-beta}
\end{equation}
Ainsi, la formule basale conserve sa fonction dérivative et la clôture
bidirectionnelle gouverne toute propagation ultérieure par caractères et tours.
Dans une carte d'unités naturelles où la coordonnée de
\(c_{\rm int}\) est fixée à un, les deux coordonnées numériques coïncident ; leurs
droites ne sont pas identifiées tacitement.

\begin{proposition}[Séparation typée de la base énergétique et du quantum d'horloge]
\label{prop:no-identificacion-base-reloj-rev8}
Écrivons \(t_0=\tau_0u_T\), avec \(\tau_0>0\), et définissons
\[
 \varepsilon_{\rm clk}
 :=\hbar_{\rm int}\frac{2\pi}{108t_0},
 \qquad
 \kappa_{\rm clk}(\tau_0)
 :=\frac{2\pi\etaRet\,10^{-34}}{108\tau_0}.
\]
Alors
\[
 \varepsilon_{\rm clk}=\kappa_{\rm clk}(\tau_0)u_E,
 \qquad
 m_0:=\frac{\varepsilon_{\rm clk}}{c_{\rm int}^2}
 =\kappa_{\rm clk}(\tau_0)\widehat c^{-2}u_M,
\]
et la section électronique en vigueur satisfait
\[
 E_e^{\rm int}
 =\frac{\mathfrak e_0}{\kappa_{\rm clk}(\tau_0)}
  \varepsilon_{\rm clk}.
\]
Par conséquent, \(E_e^{\rm int}=\mathfrak e_0\varepsilon_{\rm clk}\) si et seulement si
\(\kappa_{\rm clk}(\tau_0)=1\). Adopter \(\varepsilon_{\rm clk}\) comme base
en conservant la même coordonnée \(\mathfrak e_0\) est un nouveau choix de
repère ; cela ne constitue pas une correction de l'invariant électronique et ne découle pas
de la périodicité d'ordre 108. En particulier, la normalisation utilisée plus
loin, \(t_0=u_T\), fixe \(\tau_0=1\), non
\(\kappa_{\rm clk}=1\).
\end{proposition}

\begin{proof}
On substitue
\(\hbar_{\rm int}=\etaRet\,10^{-34}u_S\),
\(t_0=\tau_0u_T\), \(u_E=u_Su_T^{-1}\) et
\(c_{\rm int}=\widehat c,u_C\). La dernière identité résulte de la comparaison
avec \cref{eq:electron-line-valued}. Toutes les égalités ont lieu entre des
sections de la même droite ; aucune ne consulte une unité ou une masse externe.
\end{proof}

\begin{proposition}[Propagation de l'échelle de masse]
Pour chaque signature déclarée \(v\in\mathbb Z^5\),
\[
 m_v^{\rm int}:=m_e^{\rm int}\chi_{\rm mass}(v)
 \]
est une section positive de \(\mathcal L_M\), et
\[
 \frac{m_{v+w}^{\rm int}}{m_v^{\rm int}}
 =\chi_{\rm mass}(w).
\]
La construction est covariante sous changement d'échelle et ne requiert pas d'unité externe
d'énergie ou de masse.
\end{proposition}

\begin{proof}
Le caractère est adimensionnel et positif, de sorte qu'il ne modifie pas la droite de
\(m_e^{\rm int}\). L'identité des quotients est la multiplicativité de
\(\chi_{\rm mass}\). Sous un changement de repère, le numérateur et le dénominateur
acquièrent le même poids \(\mathbf M\), qui s'annule.
\end{proof}

Une carte métrologique ultérieure peut envoyer \(u_E\) sur une base d'énergie et
\(u_M\) sur la base cohérente obtenue en divisant par le carré de la vitesse
étalon. Seule cette carte autorise à exprimer la section au moyen d'un nom
d'unité externe.

\section{Relations constitutives comme isomorphismes entre droites}

Les fonctions internes satisfont
\[
 \widehat\mu=e^{\mathcal B+\mathcal E},
 \qquad
 \widehat\varepsilon=e^{-\mathcal B+\mathcal E},
 \qquad
 \widehat Z=e^{\mathcal B},
 \qquad
 \widehat c=e^{-\mathcal E}.
\]
Nous définissons les sections
\[
 \mu_{\rm int}=\widehat\mu u_\mu,
 \qquad
 \varepsilon_{\rm int}=\widehat\varepsilon u_\varepsilon.
\]

\begin{theorem}[Relations constitutives linéaires covariantes]
Dans le groupoïde des droites orientées, les identités exactes suivantes sont satisfaites
\[
 \sqrt{\frac{\mu_{\rm int}}{\varepsilon_{\rm int}}}
 =Z_{\rm int},
 \qquad
 \frac1{\sqrt{\mu_{\rm int}\varepsilon_{\rm int}}}
 =c_{\rm int}.
\]
Les quatre sections dépendent de deux caractères adimensionnels et des deux
droites \(\mathcal L_Z,\mathcal L_C\) ; elles ne constituent pas quatre échelles
indépendantes.
\end{theorem}

\begin{proof}
En coefficients,
\(\widehat\mu/\widehat\varepsilon=\widehat Z^2\) et
\(\widehat\mu\widehat\varepsilon=\widehat c^{-2}\). En dimensions,
\[
 \boldsymbol\mu-\boldsymbol\varepsilon=2\mathbf Z,
 \qquad
 \boldsymbol\mu+\boldsymbol\varepsilon=-2\mathbf C.
\]
Les racines positives sont donc des sections des droites correctes et les
identités sont covariantes.
\end{proof}

\section{Échelle de Cartan et homogénéité de l'action gravitationnelle}

La durée élémentaire \(t_0=u_T\) et la vitesse constitutive déterminent la
longueur d'une transition,
\[
 \ell_0=c_{\rm int}t_0=\widehat c\,u_L.
\]
Cette section permet d'homogénéiser la connexion affine de Cartan :
\[
 \mathcal A_{\ell_0}
 =\begin{pmatrix}
 \omega^I{}_J&\ell_0^{-1}e^I\\
 0&0
 \end{pmatrix},
 \qquad
 \mathcal F_{\ell_0}
 =\begin{pmatrix}
 F^I{}_J&\ell_0^{-1}T_{\rm tor}^I\\
 0&0
 \end{pmatrix}.
\]
Le quotient \(e/\ell_0\) est adimensionnel ; la courbure et la torsion demeurent dans
des blocs distincts du même transport.

La dimension de \(G\) n'ajoute pas un cinquième générateur, car
\(\mathbf G=-\mathbf S+5\mathbf C+2\mathbf T\). Pour tout paramètre
adimensionnel positif \(\theta_G\), nous définissons la famille covariante
\begin{equation}
 G_{\rm int}(\theta_G)
 :=\theta_G\frac{c_{\rm int}^{5}t_0^2}{\hbar_{\rm int}}
 \in\mathcal L_G^+.
 \label{eq:G-interno}
\end{equation}
Le choix \(\theta_G=1\) est la normalisation interne de Planck associée à
une transition élémentaire. Il satisfait par construction
\[
 \ell_0^2=\frac{\hbar_{\rm int}G_{\rm int}(1)}
 {c_{\rm int}^{3}}.
\]
Cette normalisation de référence n'est ni le sélecteur circulaire ni le sélecteur
holonomique de la section physique de \(G\).
Le paramètre \(\theta_G\) rend visible la liberté que ne fixent pas les
identités dimensionnelles. Une dynamique TPK produisant un caractère
gravitationnel devra le déterminer avant toute carte externe.

Soit
\[
 \kappa_{\rm int}=\frac{8\pi G_{\rm int}}{c_{\rm int}^{4}}.
\]
Dans la fonctionnelle suivante, \(M\) est une variété lisse orientée de
dimension quatre et \(\gamma\in\mathbb R\setminus\{0\}\). Si \(M\) possède un
bord, on impose des variations à support compact dans l'intérieur ou un
terme de bord avec des conditions annulant la variation au bord. Une
action de matière spinorielle présuppose, en outre, une structure de spin
compatible.
Si la cotétrade \(e^I\) prend ses valeurs dans \(\mathcal L_L\) et que la courbure
\(F^{IJ}\) est une deux-forme sans dimension après intégration de la connexion,
alors
\(\int\Sigma_{IJ}\wedge F^{IJ}\) est de dimension \(2\mathbf L\).

\begin{theorem}[Homogénéité de l'action Palatini--Cartan--Holst]
\label{thm:homogeneidad-pch}
Avec des unités explicites, la fonctionnelle gravitationnelle de dimension action est
\[
 \boxed{
 S_{\rm PCH}
 =\frac1{2\kappa_{\rm int}c_{\rm int}}
 \int\Sigma_{IJ}\wedge
 \left(\star F^{IJ}+\gamma^{-1}F^{IJ}\right)
 -\frac1{\kappa_{\rm int}c_{\rm int}}
 \int\Lambda_{\rm int}\operatorname{vol}_e
 +S_{\rm m}.}
 \]
Les deux termes géométriques appartiennent à \(\mathcal L_S\). L'écriture avec
\(1/(2\kappa)\) sans le facteur \(c^{-1}\) relève d'une convention
\(c=1\) ou possède la dimension \(\mathbf S+\mathbf C\), non la dimension action.
\end{theorem}

\begin{proof}
De \eqref{eq:dim-Gkappa},
\(\boldsymbol\kappa=-\mathbf S+\mathbf C+2\mathbf T\). En outre,
\(2\mathbf L=2\mathbf C+2\mathbf T\). Par conséquent,
\[
 -\boldsymbol\kappa-\mathbf C+2\mathbf L
 =\mathbf S.
\]
Le terme cosmologique possède la même dimension parce que
\(\boldsymbol\Lambda+4\mathbf L=2\mathbf L\). Si l'on omet
\(-\mathbf C\), le résultat est \(\mathbf S+\mathbf C\).
\end{proof}

L'équation de Cartan--Holst conserve alors son contenu variationnel
à condition que le courant de spin soit défini avec une normalisation cohérente
avec cette fonctionnelle. La réalisation dynamique doit encore fournir un
morphisme des holonomies TPK vers \((e,\omega,\Psi)\) ; la présente
section construit le codomaine dimensionnel exact et élimine l'ambiguïté des
unités dans ce morphisme.

\section{Universalité de la réalisation dimensionnelle}

Nous réunissons les résultats. Soit \(\mathcal X_{\rm HMT}\) le domaine des états
enrichis et soient des données adimensionnelles sur celui-ci :
\[
 \alpha_{\rm HMT},\quad \Adeg,\quad\Cdeg,\quad
 \widehat c,\widehat Z,\quad
 \mathcal A_*,\quad\chi_{\rm mass},\quad\mathfrak e_0,
\]
construites dans les chapitres précédents sur leurs domaines déclarés.
Considérons le domaine étendu
\[
 \mathcal X_{\rm MD}
 :=\mathcal X_{\rm HMT}\times\operatorname{Fr}_{\rm D}.
\]

\begin{theorem}[Extension dimensionnelle universelle]
\label{thm:extension-dimensional-universal}
Il existe une unique extension tensorielle, positive et
\(\mathcal G_{\rm D}\)-équivariante qui :
\begin{enumerate}
 \item réalise l'horloge d'arêtes dans \(\mathcal L_T\) ;
 \item réalise \(\widehat c\) et \(\widehat Z\) dans
       \(\mathcal L_C\) et \(\mathcal L_Z\) ;
 \item réalise l'indice déterminantal et la composante postérieure au retour
       dans \(\mathcal L_S\) ;
 \item satisfait les relations exactes entre impédance, vitesse de
       propagation, perméabilité et permittivité ;
 \item satisfait la relation interne de structure fine avec une échelle élémentaire
       de charge positive et, si un caractère \(n_Q\) est fourni, réalise
       conditionnellement \(Q(x)=n_Q(x)e_{\rm int}\) ;
 \item transporte le caractère de masse dans \(\mathcal L_M\).
\end{enumerate}
Toutes les autres droites du réseau s'obtiennent par produits tensoriels et
duaux. Deux réalisations diffèrent exactement par un unique élément de
\(\mathcal G_{\rm D}\).
\end{theorem}

\begin{proof}
Le \cref{thm:base-dimensional-hmt} identifie
\((\mathbf S,\mathbf C,\mathbf T,\mathbf Q)\) à une base libre du
réseau. Une application monoïdale sur une catégorie libre est déterminée
par sa valeur sur les générateurs. La positivité fixe les racines qui apparaissent
dans la charge et les relations constitutives. Les identités déjà démontrées garantissent que les
définitions ne dépendent pas d'une présentation alternative d'une dimension.
L'action simplement transitive de \(\mathcal G_{\rm D}\) sur les repères
prouve la dernière affirmation.
\end{proof}

La conclusion mathématique est précise. HMT fixe la généalogie des
coefficients internes et la mécanique dimensionnelle les réalise dans un réseau
ambiant minimal de rang quatre lorsque masse, longueur, durée et charge doivent
varier indépendamment. Le foncteur monoïdal des trajectoires avec signature
occupe une partie de ce réseau et n'hérite donc pas d'une affirmation de
minimalité. Les rapports et les équations homogènes sont indépendants de
toute carte. Une publication en secondes, joules, coulombs ou
électronvolts est la composition ultérieure avec un repère externe ; elle n'altère
aucun des théorèmes internes et ne peut les remplacer.
